\begin{lemma}[Terminal Kempe switch with an exact reconstruction interval]
Fix a finite loopless multigraph \(G\) with individual edge identities, a
palette of \(k\) colors, the distinguished edge \(e_0\), and an extended
Tashkinov tree \(\mathcal E\) displayed on the ordered tree \(T\), with
input coloring and generating coloring \(\varphi\). Let \(\mathcal B\)
be a hierarchy base, \(C\) a hierarchy counterexample over it, and
\(\mathit{ctx}\) a \noderef{HierarchyMinimumCounterexampleContext} for
these data. Let \(R\) be a \noderef{HierarchyReplacementPresentation}
relative to these fixed original data. Write \(\psi=R.\mathit{coloring}\),
\(H=R.\mathit{counterexample}.\mathit{hierarchy}\), \(q=H.\mathit{dividerCount}\),
and let \(T_n,D_n\) be the terminal ladder tree and raw stored-color set
of \(R\)'s replacement ETT.

Write \(|U|\) for the displayed edge count \(U.\mathit{length}\).
Assume \(|T|\geq2\), and specify the exact ordered trees
\[
 A=\noderef{HierarchyStarredStage}(C.\mathit{hierarchy},
       C.\mathit{hierarchy}.\mathit{dividerCount}),\qquad
 B=\noderef{TreeSequencePrefix}(T,|T|-1).
\]
Let \(X\) be an exact ordered prefix with
\(A\preccurlyeq X\preccurlyeq B\) and \(|A|<|X|\). Choose colors
\(a,b,\eta\) such that \(\eta\in D^H_{n,q}\),
\(a\in R.\mathit{counterexample}.\mathit{good}.\Gamma^q_\eta\), and
\[
 \noderef{HierarchyFirstMissingPosition}(\psi,T,\eta)\leq |X|.
\]
This first-missing position is totalized, with last-position fallback.
Choose a position \(j_b\) in \(B\) with
\(|A|<j_b\leq |X|\), and assume \(b\in\overline\psi(B.vertex(j_b))\).
Writing \(w=T.vertex(|T|)\), assume
\(a\in\overline\psi(w)\) or \(b\in\overline\psi(w)\).
For every displayed edge position \(t\) with \(|A|\leq t<|X|\), assume
\(\psi(T.edge(t))\ne a\); if \(\psi(T.edge(t))=b\), assume that this
individual edge is incident with \(T_n\) and that \(b\in D_n\).

Then there is a replacement presentation \(S\) relative to the same
original data, with coloring \(\tau\), such that
\[
 \overline\tau(v)=\overline\psi(v)\quad(v\in V(B)),\qquad
 \tau(T.edge(t))=\psi(T.edge(t))\quad(|A|\leq t<|X|),
\]
and both implications hold:
\[
 a\in\overline\psi(w)\Longrightarrow b\in\overline\tau(w),\qquad
 b\in\overline\psi(w)\Longrightarrow a\in\overline\tau(w).
\]
\end{lemma}
\begin{proof}
The same-good witness in \(R\), expanded using
\noderef{SameGoodHierarchy} and \noderef{SameHierarchyStages}, identifies
the replacement's last starred stage with the specified exact tree \(A\):
at zero dividers this is the common core, and at positive divider count
it is the corresponding last ordinary stage. It also identifies all
full-tail prefixes, since their definition uses only that starred tree
and the unchanged displayed tree \(T\). Since
\(|A|<|X|\leq |B|=|T|-1\), the penultimate full-tail prefix of the
replacement counterexample is exactly \(B\). Apply
\noderef{HierarchyCounterexamplePenultimateElementary} with the
properness, exactly-one-uncoloredness, lower-rung, fewer-divider,
minimality, and large-missing-set fields of \(R.\mathit{minimumContext}\).
Thus \(V(B)\) is elementary under \(\psi\).

We record first-position facts, including the fallback case. In the
finite candidate set defining \noderef{HierarchyFirstMissingPosition},
a supplied missing witness at position \(j\) bounds the minimum by \(j\).
If the returned position is strictly before \(|T|\), it is an actual
missing witness: the empty-candidate branch returns \(|T|\) and cannot
have this strict inequality. Therefore a color missing at \(B.vertex(j)\)
has first position exactly \(j\). Indeed, its first position is at most
\(j\), lies in \(B\), and is an actual missing witness; elementarity
identifies the two witness vertices, and vertex injectivity identifies
their positions. Also, any coloring agreeing with \(\psi\)'s missing
sets on \(B\) has the same totalized first positions on \(T\). If the
two positions differed, the smaller would be strictly before \(|T|\);
it would be a genuine witness in \(B\), and missing-set agreement would
contradict minimality of the larger position. This argument also covers
an empty candidate set, or a witness only at the terminal vertex.

The Gamma avoidance field of \noderef{GoodHierarchy} says that \(a\)
is missing on the last ordinary stage when \(q>0\), and on \(T_n\) when
\(q=0\). The core and nesting fields of \noderef{HierarchyBase} and
\noderef{TashkinovHierarchy} place \(T_n\preccurlyeq A\); in PE use the
ordered-union left-prefix certificate, and otherwise the core is \(T_n\).
Consequently \(a\) is missing on \(A\). Elementarity of \(B\), together
with \(j_b>|A|\), says that \(b\) is not missing on \(A\). In particular
\(a\ne b\). Let \(j_a\) be the first position missing \(a\).
Then \(j_a\leq |A|<j_b\), and the first position missing \(b\) is \(j_b\).
Both positions are genuine missing witnesses on \(B\).

Apply \noderef{ChangeLemma} to the replacement counterexample and its
minimum-context fields. Its ordering premise is \(j_a<j_b\);
its unused-color premise follows from the assumed avoidance of \(a\)
between \(A\) and \(X\), because \(j_b\leq |X|\).
Use the positive-divider conclusion when \(q>0\), and the zero-divider
conclusion with \(a\) missing on \(T_n\) when \(q=0\).
It follows that the \(a,b\)-components rooted at \(B.vertex(j_a)\)
and \(B.vertex(j_b)\) coincide. The terminal vertex \(w\) differs from
both, by exact-prefix equality and injectivity of \(T\)'s vertex map.
Now \noderef{P2905_T2_ElementaryTerminalComponent} shows that the
component \(Q\) rooted at \(w\) is disjoint from this common component,
is a two-color path component, and its swap
\(\sigma=\noderef{TwoColorComponentSwap}(G,\psi,a,b,w)\) preserves the
missing set at every vertex of \(B\).

The swap is proper by \noderef{TwoColorComponentSwapProper}.
It leaves exactly \(e_0\) uncolored: an old uncolored edge remains
uncolored, and every old colored edge either keeps its color or changes
between the two colors. Apply \noderef{StableChangeCasesOneTwo} with
the just-proved component disjointness and the same divider alternative.
Its weak-stability conclusion, expanded by \noderef{HierarchyWeaklyStable},
gives combined base stability in PE and ordinary base stability otherwise.
In either case it gives \noderef{StableColoring} on \(T_n,D_n\).

All displayed edges of \(A\) are fixed. At zero dividers this follows
from \noderef{P2813_T1_BaseDisplayedEdges} and the base stability just
obtained. At positive dividers weak stability fixes every noninitial
displayed edge of the last ordinary stage; its initial edge is \(e_0\)
and is uncolored in both colorings. An edge in the interval
\(|A|\leq t<|X|\) whose old color is \(b\) is fixed by ordinary stability,
using precisely its assumed incidence with \(T_n\) and \(b\in D_n\).
Any other edge in this interval has neither affected old color and is
fixed by the swap definition. Hence all edges of \(X\) are fixed.

We preserve the actual terminal construction route, not merely the
ordered prefix. By \noderef{P2813_T1_TerminalRoute}, in PE there is
an old augmentation sequence \(T_n\to T\); in SE or RE there is an exact
connecting edge \(f\), the connecting-edge set is literally \(\{f\}\),
and an exact seed \(T_n+f\) has an old augmentation sequence to \(T\).
The PE source is below \(A\), hence below \(X\).
The SE/RE seed has length \(|T_n|+1\leq |A|+1\leq |X|\);
it and \(X\) are prefixes of \(T\), so it is an exact prefix of \(X\).
Use \noderef{TashkinovAugmentationSequenceRestrict} for the source-to-\(X\)
route and \noderef{TashkinovAugmentationSequenceSuffix} for the
\(X\)-to-\(T\) suffix. Transport the first route by
\noderef{TashkinovAugmentationSequenceColoringTransport}, using fixed
edges of \(X\) and missing-set equality on its penultimate prefix.
This gives the required \(\sigma\)-route from \(T_n\) in PE or from
the same individual seed \(T_n+f\) in SE/RE, and retains the old
\(\psi\)-augmentation suffix from \(X\) to \(T\).

Both colors are missing on \(X\) under \(\sigma\), since their witnesses
at \(j_a,j_b\leq |X|\) are preserved on \(B\).
For any stored color \(\eta'\) whose terminal Gamma pair contains \(a\),
pairwise disjointness of Gamma pairs forces \(\eta'=\eta\); thus its
first position is at most \(|X|\) by assumption. No terminal Gamma pair
can contain \(b\), since Gamma avoidance would make \(b\) missing on
the last stage, or on \(T_n\), hence on \(A\), a contradiction.
The earlier first-position argument proves that every stored color's
totalized first position is unchanged by this swap.

All premises of \noderef{StableChangeCasesOneTwoSplitTail} now hold:
the exact \(A<X\leq B\) relations, path and divider alternatives,
interval edge equality, the actual terminal seed route, augmentation
suffix, both missing witnesses on \(X\), pointwise missing-set equality
on \(B\), unchanged stored first positions, and the affected-Gamma cutoff.
It returns an actual corresponding replacement ETT on \(T\), its own
maximum witness, a base, a good hierarchy, and a same-good witness
relative to \(R\). No equality of full segment-color sets is used here.

By \noderef{TwoColorComponentSwapMissingBehavior}, the swap exchanges
terminal missingness in both directions. The full tree remains
nonelementary: if \(a\) was missing at \(w\), the new color \(b\) is
missing both at \(w\) and at the preserved witness \(B.vertex(j_b)\);
if \(b\) was missing at \(w\), use \(a\) and \(B.vertex(j_a)\) instead.
The witnesses are distinct from \(w\) and belong to \(T\).
Together with the returned good hierarchy, maximum witness, and rung
agreement, this gives a new \noderef{HierarchyCounterexample}.
Finally apply \noderef{HierarchyReplacementPresentationTrans} to \(R\)
and these actual returned objects. It composes corresponding provenance,
same-good data, and the minimum context back to the fixed original
presentation, producing \(S\) with coloring exactly \(\sigma\).
The missing-set equality, interval edge equality, and both terminal
implications already proved therefore give all four asserted conclusions.
\end{proof}
